\documentclass[11pt,a4paper]{article}
\pdfoutput=1
% Optional starter: retain the user's existing venue and source when available.
% Copy jheppub.sty and JHEP.bst beside this file if they are not installed.
\usepackage{jheppub}
\usepackage[T1]{fontenc}
\usepackage{lmodern}
\usepackage{amsmath,amssymb,amsthm}
\usepackage{booktabs}
\newtheorem{proposition}{Proposition}[section]
\newtheorem{lemma}[proposition]{Lemma}

% Replace these neutral placeholders using information supplied by the author.
% The front matter follows the JHEP template; keep each element the author supplies.
\title{Title}
\author[a]{Author name not supplied}
\affiliation[a]{Affiliation not supplied}
% \emailAdd{...}  % one e-mail per author, when the author supplies it
\abstract{State the question, central result, reason or method, and decisive
qualification. Replace this placeholder with the actual abstract.}
\keywords{Replace with appropriate keywords}

\begin{document}
% \maketitle prints the title page and then the table of contents, as in the JHEP
% template. Keep the contents: do not add \notoc or a manual page break.
\maketitle
\flushbottom

% Replace this outline with the dependency order of the actual argument.
% A formal, computational or mixed article need not use the same sections.
\section{Introduction}
State the problem and the supported answer, and explain the route to it.

\section{Argument}
Introduce the needed objects and explain the construction, derivation or
calculation that carries the result. State assumptions where they are used.

\section{Consequences}
Interpret the result, its controlled tests and its remaining limitations.

% \acknowledgments  % only with text the author supplies
% Add appendices only for supporting material with a useful independent role.
% Supply an inspected bibliography and uncomment these lines when ready.
% \bibliographystyle{JHEP}
% \bibliography{references}
\end{document}
